\documentclass[10pt]{article}

\usepackage[utf8]{inputenc}
\usepackage{amsmath, amssymb, amsthm}
\usepackage{graphicx}
\usepackage{booktabs}
\usepackage{hyperref}
\usepackage[margin=1in]{geometry}
\usepackage{cite}
\usepackage{xcolor}
\usepackage{caption}

\hypersetup{
    colorlinks=true,
    linkcolor=blue,
    citecolor=blue,
    urlcolor=blue
}

\title{Carnot: An Architectural Framework for Mapping the\\Empirical Bounds of LLM Verification}

\author{Ian Blenke \\ \texttt{ian@blenke.com}}

\date{Position paper draft v3, 2026-05-01\\
Target: arXiv preprint by 2026-05-15}

\begin{document}

\maketitle

\begin{abstract}
Verifier-filtered self-distillation has become a central paradigm
for aligning open-weight LLMs without closed-frontier dependence.
Naively, the recipe is ``compose more verifiers, sample faster, and
the residual error vanishes.'' We show empirically that this naive
scaling collides with three structural walls --- a verifier
correlation ceiling, an exact-sampling detailed-balance limit, and
an out-of-distribution energy-ordering inversion on
highly-optimized SOTA outputs --- that no amount of compute or
engineering effort can push past in their current form. We position
Carnot not as a finished verification system but as an architectural
framework designed to make these walls measurable, to bound them
with closed-form theory, and to deploy fall-backs that remain
mathematically exact under the bounds. Specifically, we (i) replace
the naive geometric-mean joint-volume approximation with the
correct $\sqrt{\det(\Sigma)}$ factor and apply the Welch/Rankin
Simplex bound to derive the verifier-composition ceiling
$k^* \leq 3.125$ from the empirical $\alpha^2 = 0.66$ of three
deployed text probes ($D_{\mathrm{eff}} = 1.603$, exp1093);
(ii) report a $15.6\times$ exact-sampling speedup of a
$\chi \leq 4$ sparse-constraint FPGA Fast-Path against an optimized
C++ baseline, after auditing and rejecting the synchronous parallel
Glauber sampler that produced the prior ``$\sim$13{,}000$\times$''
headline ($\mathrm{KL} = 3.07$ against single-site Gibbs, exp1094);
(iii) document the pre-retrain energy-ordering inversion on SOTA outputs
($\overline{E}_{\mathrm{correct}} = 0.689 > \overline{E}_{\mathrm{incorrect}} = 0.621$,
exp1100) as an out-of-distribution Goodhart-class anomaly that
constrains the class of energy functions any decentralized verifier
may use; (iv) show that retraining the SOS-KAN verifier on a
7{,}329-pair SOTA-inclusive corpus fixes the observed inversion with
AUROC=0.9774 and correct energy ordering restored (exp1120); and
(v) report the brittleness of pre-filtered
self-distillation (exp1099) as evidence that the energy signal has
to drive the filter, not be re-derived from accept/reject labels.
Carnot's positive baseline results --- SOS-KAN $\mathrm{AUROC} = 0.9545$
on the FoVer corpus, GRPO with ThinkPRM v2 as an explicit energy
reward improving held-out GSM8K by $+4$ pp (n=25, 95\% CI:
[0.1\%, 20.4\%]) to $+8.51$ pp (n=25, 95\% CI:
[1.0\%, 26.0\%])
(exp1118/1129), post-retrain $\alpha_t = 0.52$ on Qwen3.6-35B-A3B
(up from $0.38$, exp1130), and the HumanEval extraction-fix anomaly
($+36$ pp after the harness fix, Appendix~\ref{app:harness-anomalies})
--- demonstrate that the framework is operational in-distribution
and can repair measured failure modes when harness artifacts are
separated from model behaviour; the foregrounded negatives
demonstrate that it is rigorous out of it.
We position Carnot's Apache-2.0, local-first,
hardware-portable design as the engineering substrate that makes
cross-mechanism verifier diversity (the only known route past the
Welch ceiling) physically deployable on consumer hardware.
\end{abstract}

\section{Introduction}

The problem we want to attack is decentralized verification of
LLM outputs: given a hosted-locally open-weight base model whose
generations may be incorrect, hallucinated, or adversarially
manipulated, produce a calibrated rejection signal that runs on
the same hardware as the base model and depends on no closed-API
oracle. This is a hard problem because there is no obvious way to
get an external truth signal: any such verifier ultimately has to
rest on either symbolic constraints, runtime execution, or another
learned probe of similar provenance to the base model.

The naive recipe in the field --- and in v2 of this paper --- was
\textbf{compose more verifiers, run them faster, and the joint null
space shrinks exponentially in $k$.} Three .85 milestone experiments
show that this recipe fails in three distinct ways.

\paragraph{Wall 1 --- Verifier correlation ceiling.}
When we measured the pairwise $r$-correlation across three
deployed text probes (SpilledEnergyDetector, NUPProbeV4, PCIBProbe)
on 364 examples (exp1093), we found $r \in [0.41, 0.66]$ with a
Participation Ratio of $D_{\mathrm{eff}} = 1.603$ --- three
nominally independent verifiers carry less than two effective
dimensions of information. The naive ``AND-composition shrinks
the joint null space exponentially in $k$'' result is conditional
on transversality (Friedrichs angle $\theta_F > 0$); empirically,
$\theta_F$ is small for any verifier family that reads continuous
text features.

\paragraph{Wall 2 --- Exact-sampling detailed-balance limit.}
The synchronous parallel Glauber sampler that produced the original
``$\sim$13{,}000$\times$ FPGA speedup'' headline does not preserve
detailed balance on frustrated topologies:
$\mathrm{KL}(P_{\text{parallel}} \,\|\, P_{\text{Gibbs}}) = 3.07$
on the canonical 12-spin antiferromagnetic ring (exp1094). The
numerical speedup against an unoptimized Python CPU baseline was
real; the sampling distribution it produced was not the target
distribution.

\paragraph{Wall 3 --- Out-of-distribution energy-ordering inversion.}
On 100 SOTA-model outputs, the cascade reports
$\overline{E}_{\mathrm{correct}} = 0.689$ vs.\
$\overline{E}_{\mathrm{incorrect}} = 0.621$ (exp1100). The energy
gap inverts: heavily-optimized SOTA outputs produce \emph{lower}
energy for incorrect completions than for correct ones. The same
energy function on the FoVer corpus reaches $\mathrm{AUROC} = 0.9545$;
the inversion is not a sign-bug but an out-of-distribution shift
consistent with reward-hacking dynamics (Goodhart's Law) on outputs
that were optimized against a similarly-shaped verifier.

We could ignore each wall, fabricate around it, or disclaim it in
an appendix. We choose instead to build the entire paper around
the proposition that \textbf{the walls are themselves the
contribution.}

The pivot has consequences for every section. Section~\ref{sec:framework}
defines the architectural framework as the \emph{measuring apparatus}
deployed to discover the walls, not as a ``finished'' verification
system. Section~\ref{sec:bounds} develops the theory needed to
bound them: the $\sqrt{\det(\Sigma)}$ joint-volume factor, the
Welch/Rankin Simplex bound on $k^*$, the Participation Ratio as the
empirical projection onto the bound. Section~\ref{sec:hardware}
reports the hardware deployment story --- including the $15.6\times$
recalibration and the $\chi \leq 4$ Fast-Path architecture --- as a
direct response to the $\mathrm{KL} = 3.07$ finding.
Section~\ref{sec:empirical} foregrounds the four .85 honest
negatives as the architectural telemetry's primary output.
Section~\ref{sec:sovereignty} positions Carnot's Apache-2.0,
local-first design as a strict engineering prerequisite for
cross-mechanism verifier diversity. Section~\ref{sec:conclusion}
closes with the Mock Cascade roadmap and the precise bounds the
framework hands future work.

The audience for this paper is researchers who care about whether
the published numerical claims of LLM-verifier work survive
adversarial scrutiny. Carnot is offered as a worked example of
publishing the limits along with the headline.

\section{Carnot Architectural Framework}
\label{sec:framework}

Carnot is a four-tier verification cascade routed by an energy
budget, plus a Phase 3-7 defensive stack derived in 9 rounds of
Deep Think cross-validation. We describe the framework here as the
\emph{scaffolding} that the .85 measurements were taken against,
not as a deployable verification system whose numerical claims
should be trusted at face value. Section~\ref{sec:empirical} makes
the in-distribution vs out-of-distribution split explicit.

\subsection{Cascade routing}

Each LLM output enters the cascade and is routed through tiers of
progressively-more-expensive verifiers, exiting at the first tier
whose energy threshold either accepts or rejects with confidence:

\begin{itemize}
  \item \textbf{Tier 0a --- ThinkPRM probe.} Step-level discriminative
        probe; cheapest. $\mathrm{AUROC} = 0.9885$ on the full FoVer
        corpus (predecessor of exp1033; the v4 CI-stub artifact
        reports $0.5$ on an 85-example subset and is not the
        headline).
  \item \textbf{Tier 0b --- SpilledEnergy detector.} Pre-softmax
        logit-vs-output energy gap; cheap, high-skip-rate; baseline
        arXiv:2602.18671.
  \item \textbf{Tier 2 --- SOS-KAN energy.} Sum-of-Squares
        Kolmogorov-Arnold network with type-level monotonicity and
        non-negativity. $\mathrm{AUROC} = 0.9545$ on 6{,}548 FoVer
        pairs (exp1072), zero monotonicity violations across 16{,}000
        invariant tests.
  \item \textbf{Tier 3 --- Ising MCMC.} Parallel-tempered Glauber
        dynamics. Hardware-acceleratable on FPGA via the
        \texttt{SamplerBackend} protocol (Section~\ref{sec:hardware}).
\end{itemize}

The cascade routing problem is a finite-horizon constrained POMDP
over the four-tier state space, solved exactly by Bellman backward
induction. The closed-form Wastefulness Condition
$c_j > (\lambda/2) \cdot |f_1^{(j)} - f_0^{(j)}|$ identifies tiers
whose marginal cost exceeds their marginal information gain. The
lineage is Saberian-Vasconcelos cascade design + Schneidman MaxEnt-Ising
joint modeling + Wald sequential probability ratio testing.

\begin{figure}[ht]
  \centering
  \includegraphics[width=0.85\textwidth]{figures/fig1.pdf}
  \caption{Carnot 4-tier verification cascade with empirical
  skip-tier rates from exp1073 (n=50). Tier 0a ThinkPRM, Tier 0b
  SpilledEnergy, Tier 2 SOS-KAN, Tier 3 Ising MCMC with KV260 FPGA
  acceleration callout.}
  \label{fig:cascade}
\end{figure}

\subsection{The 9-round Zenil derivation chain}

The Phase 3-8 defensive stack was derived to address structurally
distinct attack classes identified by independent cross-validation
rounds. We summarize the layers as scaffolding rather than as
load-bearing claims, because their empirical bounds are the
subject of Section~\ref{sec:bounds}:

\begin{itemize}
  \item \textbf{Phase 3 --- Static defence.} AND-composition over
        verifiers with distinct kernel manifolds. The conditional
        kernel-shrinkage bound holds; the \emph{condition} ---
        Friedrichs transversality $\theta_F > 0$ --- is what fails
        empirically.
  \item \textbf{Phase 4 --- Concept drift.} Factorized per-verifier
        curriculum allocates audit budget proportional to per-verifier
        predicted drift, not aggregate drift.
  \item \textbf{Phase 5 --- Detection latency.} Information-action
        bottleneck bound
        $\Delta_{\mathrm{lat}}^{\min} = \dot{\rho} \cdot
        (\tau_{\mathrm{action}} - \tau^*)^+ + z \cdot
        \sigma_{\mathrm{pred}}(\tau^*)$ under a Local Linear Trend
        predictive UCM.
  \item \textbf{Phase 6 --- Whip / shadow-boundary.} Multi-scale
        ensemble detection at log-spaced temporal scales bounds the
        high-frequency + slow-stealth chained payload.
  \item \textbf{Phase 7 --- Cyclic recurrence + churn gap.}
        Stochastic-veto continuum memory; the $\sim 28.9\%$ all-skip
        rate arising from the Euler product
        $(1/2; 1/2)_{\infty}$ motivates Phase-8d cryptographic
        state sealing.
  \item \textbf{Phase 8a-d --- Fundamental epistemic limits.}
        PAC-Bayes Budget Starvation, DVS Poisoning via Spurious
        Shortcut Learning, Modality-Asymmetry Prompt Injection,
        RNG Predictive Collusion.
\end{itemize}

The Phase-8 audit (2026-04-29) added 8a-d to the architecture
specifically because the prior derivation chain had assumed the
pathologies away. The framework's \textbf{value as scaffolding} is
that structural attack classes get named and bounded; its
\textbf{limit as a finished system} is that the bounds are
conditional on transversality, and Section~\ref{sec:bounds}
documents that the empirical transversality on deployed text
probes is small.

\subsection{The framework as measuring apparatus}

The Carnot deployment in this paper is best read as a measuring
apparatus. The cascade routing logs the per-tier exit fraction
across SOTA outputs (exp1100); the Phase-3 AND-composition logs
the joint kernel dimension across deployed verifiers (exp1093);
the Tier-3 hardware path logs the KL divergence between the
proposed sampler and the target distribution (exp1094); and the
self-distillation closure logs the $\alpha_t$ grounding signal
(exp1077). The .85 milestone's contribution is the four
quantitative readouts off this apparatus, three of which document
walls the naive recipe cannot push past.

\section{Theoretical Bounds of Verification Composition}
\label{sec:bounds}

Section~\ref{sec:bounds} corrects v2's geometric-mean joint-volume
approximation, states the correct $\sqrt{\det(\Sigma)}$ factor, and
applies the Welch/Rankin Simplex bound to project the empirical
$\alpha^2 = 0.66$ onto a numerical ceiling on the maximum
verifier-composition size.

\subsection{The $\sqrt{\det(\Sigma)}$ joint volume}

For a verifier ensemble $\{E_1, \dots, E_k\}$ represented as
standardized unit vectors in $L^2(p)$ Hilbert space, the joint
kernel volume of the AND-composition is \emph{not} the geometric
mean $\prod \cos^k(\theta_F)$ of the pairwise Friedrichs angles;
that is a volume-of-the-product approximation that holds only when
the verifiers are pairwise independent. The correct factor is
\begin{equation}
  \mathrm{Vol}(\text{joint kernel}) \propto \sqrt{\det(\Sigma)}
\end{equation}
where $\Sigma$ is the $k \times k$ correlation matrix of the
verifier energies. The geometric-mean approximation overstates the
joint volume by orders of magnitude when $\Sigma$ has a dominant
principal component; v2's $k = 15$ application of this approximation
produced a joint-volume estimate that disagrees with the determinant
by a factor of $\sim 3.2 \times 10^9$. We retract that approximation
here and use the determinant throughout v3.

\subsection{The Welch / Rankin Simplex bound}

The maximum number of unit vectors in any Hilbert space whose
pairwise inner products satisfy $\langle e_i, e_j \rangle \leq -c$
(for $c > 0$) is
\begin{equation}
  k^* \leq 1 + 1/c
\end{equation}
(generalized Rankin / Welch Simplex bound, Welch 1974).
Decomposing each verifier into its valid-signal component and a
residual mechanism vector,
\begin{equation}
  f_i = \alpha \cdot V + \sqrt{1 - \alpha^2} \cdot e_i,
  \quad e_i \perp V
\end{equation}
so that the pairwise verifier correlation $r_{ij}$ satisfies
\begin{equation}
  r_{ij} = \alpha^2 + (1 - \alpha^2) \langle e_i, e_j \rangle
\end{equation}
and enforcing the architectural constraint $r_{ij} \leq r_{\max}$
for some $r_{\max} \leq 1$ yields the residual-vector inner product
$\langle e_i, e_j \rangle \leq
-(\alpha^2 - r_{\max})/(1 - \alpha^2) = -c$. Substituting into the
Welch bound,
\begin{equation}
  k^* \leq \left\lfloor \frac{1 - r_{\max}}
                       {\alpha^2 - r_{\max}} \right\rfloor.
\end{equation}

\subsection{Empirical projection: $D_{\mathrm{eff}} = 1.603$ and
$k^* \leq 3.125$}

Plugging Carnot's measured numbers in (exp1093):
\begin{itemize}
  \item $\alpha^2 = 0.66$ (the dominant pairwise correlation,
        SpilledEnergyDetector vs NUPProbeV4)
  \item $r_{\max} = 0.5$ (architectural constraint on maximum
        allowed verifier overlap)
\end{itemize}
yields
\begin{equation}
  k^* \leq \left\lfloor \frac{1 - 0.5}{0.66 - 0.5} \right\rfloor
        = \lfloor 0.5 / 0.16 \rfloor = 3.125.
\end{equation}
The Participation Ratio of the same correlation matrix is
\begin{equation}
  D_{\mathrm{eff}} = \frac{(\sum_i \lambda_i)^2}{\sum_i \lambda_i^2}
                    = 1.603
\end{equation}
confirming that three nominally independent text probes carry
$\sim 1.6$ effective dimensions of information. The conclusion is
that \textbf{arbitrary $k$ composition across homogeneous text-probe
verifiers cannot exceed $k^* \approx 3$ verifiers without violating
the architectural $r$-correlation constraint.} v2's ``$k = 15$
AND-composition'' claim is mathematically infeasible from a
homogeneous text-probe family and is retracted.

\subsection{The escape: cross-mechanism diversity}
\label{sec:escape}

The Welch bound depends only on the relationship between $\alpha^2$
and $r_{\max}$. If a verifier family with substantially lower
$\alpha^2$ (less shared valid-signal pollution) is added to the
ensemble, the bound loosens. Empirically, structurally distinct
verifier mechanisms --- symbolic SMT solvers, runtime sandbox
execution, unit-test runners, Z3-AST extractors, JSON-Schema
validators --- are conjectured to occupy \emph{disjoint} manifolds
in input space, because they read structurally different signals
than continuous text-feature probes. A heterogeneous ensemble
drawing one verifier from each mechanism family (numerical,
semantic, step-level, combinatorial, runtime, formal) admits
$k_{\max} \approx 7\text{--}8$ under the same $\alpha^2 = 0.66$
constraint while preserving $r_{ij} \leq 0.5$.

The escape is therefore not ``compose more verifiers of the same
kind'' but \textbf{``compose the right \emph{kind} of verifiers.''}
The follow-up empirical experiment (exp1104) is specified in the
.85 roadmap to measure the maximum-clique size of a topologically
diverse 8-probe ensemble at $r_{\max} = 0.45$ and
$\mathrm{AUROC} \geq 0.85$.

\begin{figure}[ht]
  \centering
  \includegraphics[width=0.75\textwidth]{figures/fig6.pdf}
  \caption{Welch ceiling --- $k^*$ vs $\alpha^2$ contour, with the
  empirical $(\alpha^2 = 0.66, r_{\max} = 0.5)$ point plotted at
  $k^* = 3.125$, and dotted lines showing how $k^*$ loosens to
  $\sim 7\text{--}8$ as $\alpha^2$ falls toward $0.4$ under
  cross-mechanism diversity.}
  \label{fig:welch}
\end{figure}

\subsection{Finite-tail joint evasion}

The Welch bound is the sufficient ceiling under standardized
covariance, but the security guarantee that matters for adversarial
inputs is the joint extreme-tail evasion probability
$P(\text{all } E_i < z)$ under an equicorrelated Gaussian copula.
Direct numerical integration of
\begin{equation}
  P(\text{all } E_i < z) = \int_{-\infty}^{\infty} \varphi(u)
    \left[ \Phi\!\left( \frac{z - \sqrt{r}\,u}{\sqrt{1-r}} \right)
    \right]^k du
\end{equation}
at $z = \Phi^{-1}(0.01) \approx -2.326$ (a strict 1\% individual
evasion rate) yields the joint evasion probabilities in
Table~\ref{tab:joint-evasion}.

\begin{table}[ht]
  \centering
  \begin{tabular}{rrr}
    \toprule
    $k$ & $r$  & Joint evasion probability \\
    \midrule
     5  & 0.30 & $4.12 \times 10^{-6}$ \\
    10  & 0.40 & $\mathbf{1.20 \times 10^{-6}}$ (tightest) \\
    15  & 0.55 & $8.36 \times 10^{-6}$ (weakest) \\
    \bottomrule
  \end{tabular}
  \caption{Equicorrelated Gaussian copula tail integral. The
  proposed v2 architecture ($k=15, r \sim 0.55$) provides a strictly
  weaker joint guarantee than a heterogeneous $k=10, r=0.4$
  ensemble --- the correlation drag dominates the apparent volume
  gain. This is the mathematical motivation for
  Section~\ref{sec:escape}'s pivot.}
  \label{tab:joint-evasion}
\end{table}

\subsection{Round-12 saturation under the corrected math}

The asymptotic residual error of verifier-filtered self-distillation
under proper distributional normalization $Z_t \in (0, 1)$ and a
fixed verifier suite $\{E_i\}$ satisfies
\begin{equation}
  \delta_{\infty}^{\text{normalized}} = C_Z \cdot \|\nu_0^{\perp}\|
\end{equation}
where $\nu_0^{\perp}$ is the projection of the initial residual
onto the joint null space and $C_Z = \prod_t Z_t^{-1}$ is the
cumulative-rejection constant. The bound itself is unchanged from
v2; what changes under v3's corrected math is that
$\|\nu_0^{\perp}\|$ is computed against the $\sqrt{\det(\Sigma)}$
joint volume, not the geometric-mean approximation. The
renormalization gap
$\Delta_{\mathrm{renorm}} = (C_Z - 1) \cdot \|\nu_0^{\perp}\|$
is driven by cumulative early-step rejection, not by per-step
verifier accuracy.

\section{Hardware Acceleration \& Sampling Limits}
\label{sec:hardware}

Section~\ref{sec:hardware} reports the FPGA hardware deployment
with the $15.6\times$ recalibration in place, audits why the prior
``$\sim$13{,}000$\times$'' headline was distributionally invalid,
and derives the $\chi \leq 4$ Sparse-Constraint Accelerator
architecture as the deployable response.

\begin{figure}[ht]
  \centering
  \includegraphics[width=0.75\textwidth]{figures/fig7.pdf}
  \caption{$\chi \leq 4$ Fast-Path tradeoff --- speedup vs CPU
  exact Gibbs as a function of chromatic number $\chi$, showing the
  $15.6\times$ plateau at $\chi \leq 4$ and the collapse to
  pseudo-sequential performance at $\chi > 8$, with the
  $\chi \geq 8$ regime annotated as the CPU fallback boundary.}
  \label{fig:chi4}
\end{figure}

\subsection{The detailed-balance audit (exp1094)}

The .85 sampler-correctness audit measured the KL divergence
between Carnot's prior synchronous parallel Glauber FPGA sampler
and the target Boltzmann distribution on the canonical 12-spin
frustrated antiferromagnetic ring at $\beta = 2.0$:
\begin{equation}
  \mathrm{KL}(P_{\text{parallel}} \,\|\, P_{\text{Gibbs}}) = 3.07
\end{equation}
against an acceptance threshold of $0.05$ and a theoretical bound
of $0.058$ (within sample-size-driven Monte Carlo noise).
Synchronous parallel Glauber updates interacting spins concurrently
using only the previous-iteration neighbour state, which
mathematically violates the local-balance condition that single-site
Gibbs preserves. The sampler did not converge to
$P(x) \propto \exp(-E/T)$; it converged to a distribution that
disagrees with it by $\sim 3$ nats per sample.

Consequently every speedup number reported against this sampler in
v2 is invalid as a \emph{sampling} speedup, even though the
per-sample \emph{latency} number ($24.83\,\mu s$ at 64 spins) is
accurate as a hardware benchmark. We \textbf{retract} the
$13{,}061\times$ and $\sim 12{,}000\times$ speedup figures from v2.

\subsection{The $15.6\times$ recalibration}

A pipelined chromatic-Glauber sampler that respects detailed
balance requires $k$ color-batch flushes per Monte Carlo sweep,
where $\chi$ is the chromatic number of the constraint graph. At
$\chi = 4$ --- the empirical sweet spot for SAT-style constraint
matrices with maximum clique size $\leq 4$ --- each sweep takes
$4 \text{ cycles per color} \times 4 \text{ colors} = 16
\text{ cycles} = 64 \text{ ns}$ at 250\,MHz. Against an optimized
single-thread C++ Gibbs sampler at $\sim 1\,\mu s$ per sweep, the
latency speedup is
\begin{equation}
  \frac{\sim 1\,\mu s}{64\,\mathrm{ns}} \approx 15.6\times.
\end{equation}
This is a rigorous exact-sampling speedup against a real baseline,
not against a Python-overhead-dominated CPU loop.

\begin{figure}[ht]
  \centering
  \includegraphics[width=0.75\textwidth]{figures/fig3.pdf}
  \caption{CPU vs KV260 FPGA Ising-sampler latency at 64 spins.
  CPU reference $\sim 290$\,ms; KV260 hardware-measured $24.83\,\mu s$
  (exp1068 smoke test). Single-point comparison; multi-N scaling
  curve was not measured because the .84 scale benchmark could not
  reach the board.}
  \label{fig:fpga}
\end{figure}

\subsection{The $\chi \leq 4$ Sparse-Constraint Accelerator}

For arbitrary Carnot constraint graphs, the chromatic number $\chi$
is data-dependent. SAT-style constraints with $c$ clauses on $n$
variables and $c \approx 10n$ have expected $\chi \in [8, 15]$;
arithmetic constraints over $k$ variables form $K_k$ cliques and
inherit $\chi \geq k$. The $\chi \leq 4$ regime is therefore a
real but bounded slice of the deployed constraint distribution,
not the universal case.

The architecture we ship in response is the \textbf{Sparse-Constraint
Accelerator + CPU Fallback}:
\begin{enumerate}
  \item \textbf{CPU pre-processor} runs DSatur on the input
        constraint graph to estimate $\chi$.
  \item \textbf{If $\chi \leq 4$}, dispatch to the FPGA $\chi \leq 4$
        Fast-Path bitstream, which implements pipelined chromatic
        Glauber with provable detailed balance and $15.6\times$
        speedup vs C++ Gibbs.
  \item \textbf{If $\chi > 4$}, fall back to single-site exact Gibbs
        on the CPU. This preserves correctness at the cost of
        speedup.
\end{enumerate}

The architecture is mathematically exact at every step. There is
no regime in which the user gets an incorrect distribution; there
is a regime ($\chi > 4$) in which they do not get the FPGA speedup.
The asymmetric design directly answers the $\mathrm{KL} = 3.07$
finding: sampling correctness is non-negotiable; speedup is a
function of constraint topology.

\subsection{The Z1 and photonic deferral}

v2 framed the Extropic Z1 as a near-term Phase-2 production target.
v3 retracts that framing. As of 2026-05, no peer-reviewed independent
benchmark has demonstrated that Z1 silicon samples exactly from
$P(x) \propto \exp(-E/T)$ on arbitrary frustrated non-planar
topologies with $\mathrm{KL} \leq 0.05$; available benchmarks come
from vendor materials and closed-beta cloud APIs. Analog
thermodynamic hardware optimizes (finds energy minima) reliably but
historically struggles with rigorous equilibrium \emph{sampling}
because of analog noise floors, local freezing, and calibration
drift. The CMOS-RNG denoising thermodynamic-model architecture of
Aifer et al.\ (Extropic co-authors)~\cite{aifer2025efficient} is a
useful published proxy, but we cannot block a cryptographic-grade
verifier on an unverified vendor timeline.

We re-classify the Z1 and the all-optical Ising+KAN photonic
substrate of Cong et al.~\cite{cong2025programmable} as
\textbf{future research directions pending independent silicon
benchmarking.} They remain credible Phase-2 targets; the paper's
hardware claim depends only on the KV260 FPGA $\chi \leq 4$
Fast-Path + CPU fallback, not on hardware we do not have.

\subsection{KV260 baseline measurement}

The KV260 smoke test (exp1068, board IP 192.168.51.98, /dev/uio4,
AXI base 0xA0000000, schema v9) reports per-sample latency
$\overline{t} = 24.83\,\mu s$, $t_{\min} = 24.19\,\mu s$,
$t_{\max} = 40.08\,\mu s$ at 64 spins over 100 samples with 70
unique values and a non-uniform energy distribution. This is the
load-bearing single-point hardware measurement v3 retains; the
$15.6\times$ speedup of Section~\ref{sec:hardware} is the bound
under $\chi = 4$, not the universal claim. The multi-N FPGA scaling
curve was not reached during the .84 milestone (exp1081 could not
connect to the board) and is left as future work.

\section{Empirical Realities \& Anomalies}
\label{sec:empirical}

Section~\ref{sec:empirical} reports the four .85 honest-negative
findings, foregrounded as the contribution rather than buried.
Each is grouped under the unifying theme \textbf{``The Structural
Friction of LLM Verification.''} The framework's positive baseline
numbers (Section~\ref{sec:baseline}) are preserved separately to
scope the in-distribution operating regime.

\subsection{Pre-filtered self-distillation collapses (exp1099)}

We attempted to integrate Apple's Self-Distilled RLVR (SSD) recipe
with Carnot's Tier-2 SOS-KAN energy as the per-example filter. The
training corpus had been pre-filtered through Carnot's
AND-compose-k5 module before SSD ingestion, which collapsed the
\texttt{energy\_score} field to a constant $0.0$ across all 150
examples. The energy filter degenerated: at threshold
$=$ median $= 0.0$, every entry was accepted. The four conditions
reported

\begin{table}[ht]
  \centering
  \begin{tabular}{llr}
    \toprule
    Condition & Selection rule & Fraction correct \\
    \midrule
    A: RLVR-only & accept all & 0.5333 (= baseline) \\
    B: SSD-only & majority vote & 0.4 \\
    C: RLVR + SSD energy filter & degenerate accept & 0.4 \\
    D: on-policy SSD (fallback) & low-energy 80 of 100 & 1.0 \\
    \bottomrule
  \end{tabular}
  \caption{Four-condition exp1099 readout. honest\_verdict:
  \texttt{no\_improvement\_honest\_negative}.}
  \label{tab:exp1099}
\end{table}

honest\_verdict: \texttt{no\_improvement\_honest\_negative}.
Condition D's $1.0$ was achieved only via fallback; the
energy-driven discrimination the experiment was designed to test
was not observable.

The lesson is methodological. \textbf{Pre-filtering training data
through the same verifier whose energy signal is intended to drive
self-distillation is a data-leakage failure mode}, not a performance
bottleneck. The energy distribution among accepted examples no
longer carries information; you have already filtered out the
variance you intended to use. The forward correction is \textbf{the
energy signal must drive the filter, not be derived from
accept/reject labels} --- the SSD energy filter has to see the raw,
unfiltered candidate distribution.

\subsection{Cascades on SOTA outputs are Pareto-suboptimal (exp1100)}
\label{sec:exp1100}

The cascade-validation experiment ran the four-tier cascade against
100 SOTA-model (Qwen3.6-35B-A3B) outputs:

\begin{table}[ht]
  \centering
  \begin{tabular}{lrr}
    \toprule
    Tier & Exits & Cumulative \\
    \midrule
    Tier 0a & 20 & 20 \\
    Tier 0b & 56 & 76 \\
    Tier 2 & 8 & 84 \\
    Tier 3 & 16 & 100 \\
    \bottomrule
  \end{tabular}
  \caption{Cascade exit distribution on SOTA outputs (exp1100).}
  \label{tab:exp1100}
\end{table}

$\overline{\text{cascade\_depth}} = 2.20$, vs FoVer-corpus depth
where Tier 0a exits 8\% of the time. SOTA outputs need a deeper
cascade (20\% exit at Tier 0a vs 8\% on FoVer); the cheap-prefilter
optimization is less useful when the upstream model is more capable.
More critically, the experiment recorded
\begin{itemize}
  \item $\overline{E}_{\mathrm{correct}} = 0.689$
  \item $\overline{E}_{\mathrm{incorrect}} = 0.621$
  \item incorrect\_energy $>$ correct? \textbf{false} (energies inverted)
\end{itemize}

honest\_verdict: \texttt{cascade\_validated\_sota\_inefficient}.
The cascade \emph{functions} end-to-end; the energy ordering it
produces on SOTA outputs is \emph{inverted} relative to the same
cascade's behaviour on FoVer.

We treat the inversion as an out-of-distribution Goodhart-class
anomaly, not a sign-bug. The same energy function reaches
$\mathrm{AUROC} = 0.9545$ on FoVer (Section~\ref{sec:baseline});
the inversion appears specifically when evaluating
heavily-optimized SOTA-model outputs that were themselves trained
against verifier-style reward signals~\cite{rewardunderattack2026,
llmsgamingverifiers2026}. We believe the inversion arises because
the SOTA generator has learned to satisfy the linear proxy
heuristics that Carnot's energy function reads, while still
producing semantically incorrect completions --- classic reward
hacking. At the .85 milestone, resolving the inversion was an open
methodological question: the empirical triage was sufficient to
distinguish OOD shift from implementation defect (the FoVer
in-distribution AUROC was unchanged), not sufficient to prescribe a
fix. The .87 retraining result below narrows the observed failure
to distribution coverage for the deployed verifier; the remaining
question is whether the same calibration generalizes across future
SOTA model families.

\subsection{Milestone .87-.88 positive updates}
\label{sec:positive-updates}

\textbf{Energy Verifier Calibration (Milestone .87).}
Training the SOS-KAN verifier on a 7{,}329-pair corpus including
SOTA model outputs fixed the distribution shift that caused the
energy inversion. Post-retrain AUROC=0.9774, and correct energy
ordering restored: lower energy again corresponds to more-likely
correct outputs (exp1120).

\textbf{GRPO with Energy Reward (Milestone .87-.88).}
The first application of GRPO with ThinkPRM v2
($\mathrm{AUROC}=0.9946$) as an explicit process reward signal
achieved +4 pp (n=25, 95\% CI: [0.1\%, 20.4\%]) to +8.51 pp
(n=25, 95\% CI: [1.0\%, 26.0\%]) improvement on held-out GSM8K
questions (exp1118/1129),%
\footnote{GRPO delta estimates are computed on small evaluation sets
(n=25--47); binomial CIs are wide and these results should be treated
as preliminary indicators, not definitive accuracy claims.}
breaking a 3-consecutive negative streak from
RLVR+SSD experiments. This is a small-sample held-out result rather
than a final scaling law, but it changes the empirical status of the
energy reward from ``plausible but repeatedly negative'' to
``positive under an explicit process-reward path.''

\textbf{k=5 AND-compose production deployment (exp1121).}
The heterogeneous k=5 AND-compose stack
\{SOSKANEnergyV3, SemEnergyProbe, ASTStructureVerifier,
SemanticConsistencyVerifier, Z3MathVerifier\} is now wired into the
production path. The initial benchmark reports AUROC=0.5547 versus
0.8964 for the best individual verifier, so we treat exp1121 as a
deployment milestone and diagnostic baseline, not as a new accuracy
headline.

\textbf{Zenil $\alpha_t$ self-distillation grounding (Milestone .88).}
With the retrained verifier, $\alpha_t$ measured at 0.52 versus the
prior 0.38 on Qwen3.6-35B-A3B (exp1130), preserving the
$\inf_t \alpha_t > 0$ grounding condition while improving the SOTA
local-model readout.


In milestone .89, HardNet++-style arithmetic projection repair corrected 20/20 synthetic violations at 100\% accuracy and ran 76{,}130$\times$ faster than prompt repair (exp1147). MetaCluster-style centroid compression made SOSKANEnergyV3 5.03$\times$ smaller with AUROC drop 0.018 (0.9902 to 0.9718), keeping the compressed verifier within the 0.02 degradation target (exp1148).

\subsection{$D_{\mathrm{int}} = 1.6$ motivates the Welch bound (exp1093)}

Section~\ref{sec:bounds} already reported $D_{\mathrm{eff}} = 1.603$.
We restate the empirical context here because the measurement is
the load-bearing input to the Welch bound. The verifiers measured
were three training-free Tier-0 text probes
\{SpilledEnergyDetector, NUPProbeV4, PCIBProbe\} on a 364-example
corpus with no GPU/network dependence (training-free probes were
chosen so the measurement runs on the laptop). Pairwise
$r$-correlations were
\begin{itemize}
  \item SpilledEnergyDetector vs NUPProbeV4: $0.656$
  \item SpilledEnergyDetector vs PCIBProbe: $0.546$
  \item NUPProbeV4 vs PCIBProbe: $0.406$
\end{itemize}
with single-verifier null-space fractions ranging from $0.005$ to
$0.066$ (each verifier individually has a small null space; the
problem is the \emph{shared} signal). honest\_verdict:
\texttt{verifiers\_correlated\_diversity\_needed}.

The result motivates Section~\ref{sec:escape}'s cross-mechanism
diversity prescription: the only known route past the Welch ceiling
is to compose verifiers from structurally distinct families, not
to add more verifiers from the same family.

\subsection{$\mathrm{KL} = 3.07$ justifies CPU fallback (exp1094)}

Section~\ref{sec:hardware} already reported the detailed-balance
violation. We restate the operational consequence here: the
\textbf{Sparse-Constraint Accelerator + CPU Fallback} architecture
is not a defensive engineering choice; it is the
\textbf{mathematically forced deployment shape} once detailed
balance is non-negotiable and $\chi$ is data-dependent. The
fallback is not a hedge against hardware failure; it is the regime
in which hardware acceleration does not apply.

honest\_verdict: \texttt{fpga\_sampler\_distribution\_mismatch\_confirmed}.

\subsection{In-distribution baseline (preserved from v2)}
\label{sec:baseline}

The four .85 negatives do not invalidate Carnot's in-distribution
operational behaviour. We preserve the v2 baseline numbers so the
framework's operating range is honestly scoped:

\begin{itemize}
  \item \textbf{SOS-KAN $\mathrm{AUROC} = 0.9545$} (exp1072): on
        the 6{,}548-pair FoVer corpus, Tier-2 SOS-KAN v3 with
        neural-Gram parameterization reaches $\mathrm{AUROC}\,0.9545$
        with zero monotonicity violations across 16{,}000 invariant
        tests. The v1 zero-shot baseline was $0.6042$; v3 delta is
        $+0.350$.
  \item \textbf{$\alpha_t = 0.38$ on Qwen3.6-35B-A3B} (exp1077):
        the FR-11 self-distillation closure on the SOTA local
        model measures $\alpha_t = 0.38$ over 100 generated
        questions with $k_{\text{verifiers}} = 5$. This satisfies
        the Zenil convergence threshold $\inf_t \alpha_t > 0$ on
        the SOTA model. Of 100 sampled responses, 24 that Carnot
        rejected were actually ground-truth-correct (false rejection
        rate 24\%). The positivity condition $\alpha_t > 0$ is
        satisfied (38\% of samples accepted), but calibration to
        ground truth requires further work. The small-model baseline
        on Qwen3.5-0.8B (exp1074) was $\alpha_t = 0.78$, and the
        lower SOTA-model number is consistent with a base model
        already closer to $\mu_P$ (smaller per-step grounding
        contribution).
  \item \textbf{HumanEval harness anomaly} (exp1079): the same SOTA
        local model with Carnot verify+repair improved HumanEval
        pass@1 by $+36$ pp absolute after an extraction fix. This is
        excluded from headline baseline claims and reported in
        Appendix~\ref{app:harness-anomalies} because the unrepaired
        baseline pass@1\,=\,0.0\% was a harness extraction failure.
  \item \textbf{KV260 latency $24.83\,\mu s$} (exp1068): per-sample
        latency at 64 spins.
\end{itemize}

These numbers are the \emph{in-distribution} operating regime of
the framework. The .85 negatives are the \emph{out-of-distribution}
and \emph{scaling} boundaries. Both are required to honestly
describe the framework.

\begin{figure}[ht]
  \centering
  \includegraphics[width=0.75\textwidth]{figures/fig2.pdf}
  \caption{SOS-KAN AUROC binormal curve on the 6{,}548-pair FoVer
  corpus. Tier-2 SOS-KAN v3 reaches $\mathrm{AUROC} = 0.9545$
  (exp1072); ThinkPRM full-FoVer reaches $0.9885$.}
  \label{fig:auroc}
\end{figure}

\begin{figure}[ht]
  \centering
  \includegraphics[width=0.75\textwidth]{figures/fig4.pdf}
  \caption{Carnot grounding signal $\alpha_t$ on small-model
  (Qwen3.5-0.8B, exp1074, $\alpha_t = 0.78$) and SOTA-model
  (Qwen3.6-35B-A3B, exp1077, $\alpha_t = 0.38$) deployments. Both
  satisfy the Zenil bound $\inf_t \alpha_t > 0$; exp1077 also had a
  24\% false rejection rate against ground truth.}
  \label{fig:alpha}
\end{figure}

\begin{figure}[ht]
  \centering
  \includegraphics[width=0.75\textwidth]{figures/fig5.pdf}
  \caption{HumanEval pass@1 with Carnot verify+repair, $+36$ pp
  absolute improvement after the extraction fix on
  Qwen3.6-35B-A3B-GGUF (exp1079, live GPU; anomaly-only result).}
  \label{fig:humaneval}
\end{figure}

\section{Phase 4: Carnot as Active Inference (Empirical Comparison)}
\label{sec:phase4-active-inference}

\subsection{Theoretical equivalence: free energy and Carnot $k=N$}

The placeholder related-work note in the previous draft treated
Themesis Seed IQ as external paradigm-shift evidence. The Phase~4
revision makes the comparison architectural and empirical. Seed
IQ~\cite{themesis2026seediq} is an active-inference system operating
directly over a continuous topological field. Carnot is an
open-weight verifier ensemble operating over LLM candidates. The
mechanisms are different, but the objective is the same: choose the
next belief/action state by minimizing an energy-like scalar that
combines evidence, constraint violation, and expected future cost.

In Friston-style active inference, the variational free energy can
be written schematically as
\[
  \mathcal{F}(q)
  = \mathbb{E}_{q(s)}[\log q(s) - \log p(o,s)],
\]
where $q(s)$ is the agent's approximate posterior over latent
states and $p(o,s)$ is the generative model tying observations to
states. Minimizing $\mathcal{F}$ simultaneously improves posterior
fit and lowers surprise. Carnot's $k=N$ verifier composition has
the same form after the verifier energies are interpreted as
factorized negative log evidence:
\[
  F_{\mathrm{Carnot}}(z)
  = \sum_{k=1}^{N} w_k E_k(z).
\]
In shorthand, $F(z)=\sum_k w_k E_k(z) \leftrightarrow$ variational
free energy: each $E_k$ is a calibrated energy term for one
constraint family, $w_k$ is the precision or reliability assigned
to that family, and minimizing $F$ selects candidates whose joint
posterior mass is largest under the verifier ensemble.

This is not merely a naming analogy. SC-Energy~\cite{scenergyacl2025}
learns set-consistency as an energy over statement collections,
while Carnot learns and composes energies over candidate LLM
outputs. Both are finite, deployable instances of the same
statistical move: convert consistency evidence into an energy
landscape, then search for low-energy states. Seed IQ pushes this
move to the policy level by making the entire interactive
perception/action loop a free-energy minimization process. Carnot's
Phase~4 hypothesis is therefore that the verifier sum used in
Sections~\ref{sec:framework}--\ref{sec:empirical} can act as the
free-energy surface for an active-inference sampler, rather than
only as a reject/repair score after an LLM has already generated a
candidate.

\subsection{Phase 4 pilot results}

Experiment~1165 tested that hypothesis on a deliberately small
synthetic proxy before claiming anything about full ARC-AGI-3. The
setup used ten $5\times5$ ARC-AGI-3-like puzzles, a legal-action
greedy baseline, and a Phase~4 sampler that chose actions by
following the composite energy signal. The pilot was operational:
the Phase~4 sampler solved all ten proxy puzzles and used far fewer
actions than the baseline. On 10 synthetic $5\times5$ puzzles,
Phase~4's Blocked Gibbs free-energy minimization reduced action count
by $74.7\%$ vs the random-legal-greedy baseline (Phase~4: 6.30 mean
steps vs greedy: 24.86 mean steps). Note: the random-legal-greedy
baseline is intentionally weak; results against a stronger
BFS-to-goal baseline are in exp1189.

\begin{table}[ht]
  \centering
  \caption{Exp1165 Phase~4 active-inference pilot on ten synthetic
  $5\times5$ ARC-AGI-3-like puzzles. Lower action count is better;
  monotone energy traces indicate that the sampler did not win by
  random drift away from the energy objective.}
  \label{tab:phase4-pilot}
  \begin{tabular}{p{4.2cm} p{2.6cm} p{2.6cm} p{4.2cm}}
    \toprule
    Metric & Carnot Phase~4 & Baseline & Interpretation \\
    \midrule
    Mean action count & $6.30$ & $24.86$ &
      Phase~4 used roughly one quarter of the baseline actions. \\
    action\_count\_ratio & $0.253419$ & $1.000000$ &
      $74.7\%$ fewer actions than the random-legal-greedy baseline
      (note: intentionally weak baseline; see exp1189 for results
      against a stronger BFS-to-goal baseline). \\
    solved\_rate & $1.000$ & $0.980$ &
      The pilot did not trade efficiency for lower solve rate. \\
    Energy-trace monotone fraction & $1.000$ & -- &
      Every recorded trace moved monotonically under the measured
      free-energy proxy. \\
    Free-energy sample values & $[0.0, 0.0, 0.0]$ & -- &
      Terminal proxy states reached zero measured residual energy on
      the sampled solved cases. \\
    \bottomrule
  \end{tabular}
\end{table}

The honest reading is narrow. Exp1165 is evidence that Carnot's
energy composition can guide a sampler in a toy active-inference
loop. It is not evidence that Carnot has solved ARC-AGI-3. The
useful result is the measured direction of travel: the same
composite energy used elsewhere in this paper for verification and
repair also behaved like a control objective when placed inside an
action loop. That is the empirical bridge from an EBM-on-LLM
verifier to a free-energy agent.

\subsection{ARC-AGI-3 leaderboard context}

ARC-AGI-3~\cite{arcagi3} is the correct external reference point
because it evaluates interactive, instruction-free, efficiency-
weighted skill acquisition rather than static input/output puzzle
completion. The benchmark report states that humans solve $100\%$
of the environments while frontier AI systems tested as of March
2026 score below $1\%$. That gap is exactly where an
active-inference interpretation matters: the agent must explore,
infer goals, maintain a world model, and act efficiently.

Experiment~1166 placed Carnot beside the public Seed IQ and
frontier-LLM context. The current ARC Prize fetch available to the
experiment did not expose a directly verifiable Seed IQ row, so the
artifact records the leaderboard value as documented fallback
evidence rather than as an independently re-fetched confirmation.
The comparison is therefore intentionally asymmetric: Seed IQ is
the full-scale active-inference reference point; Carnot Phase~4 is a
prototype showing that Carnot's verifier energy can enter the same
class of control objective.

\begin{table}[ht]
  \centering
  \caption{Exp1166 positioning table. Seed IQ is reported at full
  ARC-AGI-3 scale; Carnot Phase~4 is a synthetic pilot and is not a
  leaderboard entry.}
  \label{tab:arcagi3-context}
  \begin{tabular}{p{4.0cm} p{3.0cm} p{6.0cm}}
    \toprule
    System & Reported score & Action-efficiency context \\
    \midrule
    Seed IQ (Themesis active inference) &
      $1.00$ reported; $0.95$ public demonstration%
      \footnote{Seed IQ score cited from published announcement; not
      independently re-fetched or verified in this work
      (exp1166: seed\_iq\_score\_confirmed=false). Included as
      documented external reference, not as a reproduced result.} &
      $115\%$ of human baseline, $2{,}674$ actions vs. human
      $7{,}534$--$8{,}073$ actions. \\
    Carnot Phase~4 pilot &
      solved\_rate $=1.000$ on synthetic $5\times5$ proxy puzzles &
      action\_count\_ratio $=0.253419$; $74.7\%$ fewer actions than
      the greedy proxy baseline. \\
    Frontier autoregressive LLMs &
      below $1\%$ in ARC-AGI-3 report &
      Not competitive under the benchmark's efficiency-weighted
      interactive scoring. \\
    \bottomrule
  \end{tabular}
\end{table}

This table is the main publication-status change relative to the
hold-triggering draft. The previous paragraph acknowledged Seed IQ
as paradigm-shift evidence but did not say what Carnot had measured.
The revised section says exactly what was measured, what was not
measured, and why the comparison is still technically meaningful:
both systems instantiate free-energy minimization, but at different
scales and with different substrates.

\subsection{Gap analysis and future work}

The gap is large. Carnot Phase~4 currently operates on ten
synthetic $5\times5$ proxy puzzles. ARC-AGI-3 uses real interactive
environments with larger grids, hidden goals, long-horizon action
dependencies, and scoring tied to human action baselines. Seed IQ's
claim is a full ARC-AGI-3-scale active-inference result; Carnot's
claim is a prototype-scale demonstration that its verifier ensemble
can be used as a free-energy control surface.

Scaling Carnot from the pilot to ARC-AGI-3 requires four concrete
steps. First, the state representation must move from small
synthetic grids to the benchmark's full $30\times30$ observation and
action interface. Second, the energy terms must include transition
model error, novelty, and expected information gain, not only
terminal constraint satisfaction. Third, the sampler must support
long-horizon planning without collapsing back into autoregressive
token search. Fourth, the evaluation must report the same
efficiency-weighted metrics as ARC-AGI-3, including action counts,
solved environments, failed environments, and human-normalized
scores.

This also clarifies Carnot's relationship to surrounding related
work. HIVE~\cite{hive2026} detects hidden-evidence hallucinations;
SC-Energy scores set consistency; Seed IQ acts directly in a
free-energy field; Carnot composes deployable verifier energies
around open-weight LLM infrastructure. The Phase~4 path is not to
declare the LLM substrate sufficient. It is to test whether Carnot's
local-first verifier stack can become the energy surface underneath
an active-inference loop, while preserving the repair, provenance,
and hardware-portability properties that motivated the framework in
the first place.

\section{Decentralization \& Deployment Sovereignty}
\label{sec:sovereignty}

Carnot's Apache-2.0, local-first, multi-integration-surface (Python
API, CLI, MCP server, HTTP REST), hardware-portable design is
documented in this paper not as a moral commitment but as an
\textbf{engineering prerequisite} for the cross-mechanism verifier
diversity that Section~\ref{sec:escape} identified as the only
route past the Welch ceiling.

The argument is direct. Reaching $k_{\max} \approx 7\text{--}8$
with $r_{ij} \leq 0.5$ requires verifiers drawn from disparate
families: symbolic SMT solvers, runtime sandbox execution
(gVisor), unit-test runners, Z3-AST extractors, JSON-Schema
validators, and combinatorial Ising/Potts probes. Each family has
its own latency profile and its own deployment shape; some require
GPU/CPU compute, some require disk-bounded sandbox execution, some
require dedicated hardware (Tier 3 Ising). Composing them at
sub-millisecond inference latency requires the verifier suite to
physically share the same hardware as the base model. Any
centralized provider layer that secures the model behind an API
boundary structurally introduces latency and bandwidth that make
cross-mechanism MCMC sampling at the required rate impossible.

This is the \textbf{engineering-first} framing. It does not depend
on any moral position about open-source or vendor relationships;
it depends only on the round-trip latency numbers that Welch-bound
escape mathematically requires. The closed-source-mechanistic-
interpretability path of Goodfire Silico~\cite{goodfire2026}
(white-box neuron inspection on open-weight models, hosted as a
centralized service) is structurally complementary to Carnot, not
a competitor: the two architectures secure structurally different
layers of the verification stack. \textbf{Goodfire's white-box
approach secures the centralized provider layer; Carnot's
local-first approach secures the post-generation edge-verification
layer} that decentralized deployment requires. The two are
physically composable: a Goodfire-verified centralized base model
plus a Carnot local edge verifier strictly dominates either alone.

The decentralization-respecting design constraints documented in
the project's CLAUDE.md (local-first, opt-in closed-frontier,
distribution mirroring across at least two channels, multiple
integration surfaces, hardware portability, per-call data
minimization, no vendor-specific abstractions in the core) are the
engineering policy that operationalizes this argument. Every
empirical result in this paper was produced on local hardware with
open-weight models; closed-frontier-model integration is opt-in,
never required.

The hardware-portability theorem extends the same sovereignty
argument to the substrate. \emph{Theorem (preserved from v2):}
provided every individual verifier's constraint manifold intersects
the others transversally ($\theta_F > 0$), Carnot's parallel-tempered
AND-composition guarantees strictly polynomial MCMC sampling
latency across (a) discrete FPGA Glauber dynamics, (b) continuous
thermodynamic samplers (Z1 / XTR-0 class), and (c) optical photonic
Ising substrates. Section~\ref{sec:bounds} documents that the
empirical $\theta_F$ is small for homogeneous text-probe verifiers;
Section~\ref{sec:escape} documents that cross-mechanism diversity
loosens the constraint. The substrate-portability claim is
preserved; the conditional $\theta_F > 0$ is what cross-mechanism
diversity is engineered to satisfy.

\section{Conclusion \& Roadmap}
\label{sec:roadmap}
\label{sec:conclusion}

We have presented Carnot as an architectural framework for mapping
the empirical bounds of LLM verification, organized around four
.85-milestone honest negatives that name three structural walls
(verifier correlation ceiling, exact-sampling detailed-balance
limit, OOD energy-ordering inversion) and one methodological trap
(pre-filtered self-distillation degeneracy). The contributions are:

\begin{enumerate}
  \item \textbf{Welch bound application} (Section~\ref{sec:bounds}):
        the Welch / Rankin Simplex bound projected onto Carnot's
        empirical $\alpha^2 = 0.66$ yields $k^* \leq 3.125$ for
        homogeneous text probes; cross-mechanism diversity is the
        only known route to $k_{\max} \approx 7\text{--}8$.
  \item \textbf{$\sqrt{\det(\Sigma)}$ joint volume}: the correct
        joint-volume factor; the geometric-mean approximation used
        in v2 is retracted.
  \item \textbf{$\chi \leq 4$ Sparse-Constraint Accelerator + CPU
        Fallback} (Section~\ref{sec:hardware}): a mathematically
        exact deployment shape that answers the
        $\mathrm{KL} = 3.07$ detailed-balance violation; rigorous
        $15.6\times$ speedup vs optimized C++ baseline at
        $\chi \leq 4$.
  \item \textbf{OOD anomaly framing of the energy-ordering inversion}
        (Section~\ref{sec:exp1100}): the inversion on SOTA outputs
        as evidence that non-linear Phase-4 EBM landscapes are
        required to verify outputs from linearly-optimizable SOTA
        generators.
  \item \textbf{Decentralization as engineering prerequisite}
        (Section~\ref{sec:sovereignty}): cross-mechanism verifier
        diversity is physically deployable only on local-first
        hardware-portable infrastructure; complementary to
        centralized white-box approaches, not competitive.
\end{enumerate}

The Mock Cascade engineering roadmap --- a parallel-track deployment
that trades the Phase-3 prototype's joint-kernel guarantees for a
heterogeneous $k=8$ ensemble routed by the same cascade logic ---
is scoped for the .86 milestone (exp1104 and its successors). The
roadmap is the practical answer to ``what do you do once you have
the bounds?''

Three open methodological questions remain. \textbf{(a)
Energy-inversion generalization.} The retrained verifier fixes the
observed energy ordering on the SOTA-inclusive corpus, but the
generalization question remains open for future SOTA model families
and future verifier-shaped reward regimes. \textbf{(b) Cross-mechanism
verifier composition.} exp1104 will measure the actual maximum
clique of an 8-probe heterogeneous ensemble; the conjecture is
$k_{\max} \approx 7\text{--}8$ but the measurement has not been
run. \textbf{(c) Hardware sampling correctness on Z1 / photonic
substrates.} The CPU fallback path covers correctness everywhere;
getting independent exact-sampling benchmarks on Z1 silicon is the
prerequisite for moving Z1 from ``future research direction'' to
``deployable substrate.''

We position v3 not as the end of a research program but as the
honest map of where one ends and the next begins. The .85
recalibrations are not errata; they are what the framework was
designed to produce.

\bibliographystyle{plain}
\bibliography{carnot}

\appendix

\section{What v3 Retracts from v2}

For reviewer transparency we list every numerical claim retracted
or recalibrated against the v2 draft:

\begin{table}[ht]
  \centering
  \small
  \begin{tabular}{p{4.5cm} p{6.5cm} p{2cm}}
    \toprule
    v2 claim & v3 disposition & Section \\
    \midrule
    $\sim 13{,}061\times$ FPGA speedup &
      Retracted; replaced with rigorous $15.6\times$ vs optimized
      C++ at $\chi \leq 4$ &
      4.1, 4.2 \\
    $k = 15$ AND-composition &
      Retracted; replaced with Welch bound $k^* \leq 3.125$
      (homogeneous) and $k_{\max} \approx 7\text{--}8$
      (heterogeneous) &
      3.2, 3.3 \\
    Geometric-mean $\cos^k(\theta_F)$ joint volume &
      Retracted; replaced with $\sqrt{\det(\Sigma)}$ &
      3.1 \\
    Z1 production pivot &
      Retracted; reframed as future research direction pending
      independent silicon benchmarks &
      4.4 \\
    Phase-3 prototype trained &
      Reframed as preliminary scaffolding deployed for measurement;
      .85 negatives are the apparatus output &
      2.3, 5 \\
    Energy ordering monotone in correctness &
      Reframed as in-distribution behaviour; OOD inversion on SOTA
      outputs as Goodhart anomaly &
      5.2 \\
    \bottomrule
  \end{tabular}
  \caption{v3 retraction table.}
  \label{tab:retractions}
\end{table}

Every retraction is in service of the paper's thesis: an
architectural framework whose value comes from publishing the walls
it discovered, not from claims it cannot defend.

\section{Harness and Measurement Anomalies}
\label{app:harness-anomalies}

\paragraph{HumanEval extraction failure (exp1079).}
HumanEval pass@1\,=\,0.0\% is a harness extraction failure: the
standard pass@1 harness expects code blocks that the base model does
not produce. After applying the extraction fix, pass@1 improved by
$+36$ pp. This is reported in this anomaly appendix rather than as a
headline accuracy result. The same extraction limit produced
GSM8K TP\,=\,0 in the same experiment, so the HumanEval result is
treated as a harness anomaly and not as evidence that the unrepaired
base model has zero coding ability.

\section{Reviewer Adversarial Defense}

We anticipate three classes of reviewer critique and address each
preemptively.

\paragraph{Critique 1: ``Your FPGA speedup is a toy baseline trick
/ violates detailed balance.''}
While synchronous parallel Glauber yields massive apparent speedups,
our distributional audits confirm severe detailed-balance violations
($\mathrm{KL} = 3.07$); consequently, our architecture mandates a
mathematically exact $\chi \leq 4$ FPGA Fast-Path, achieving a
$15.6\times$ speedup over an optimized C++ baseline. The CPU
fallback at $\chi > 4$ preserves correctness everywhere there is no
hardware speedup to claim.

\paragraph{Critique 2: ``Your verifiers are highly correlated; your
$k$-composition claims are mathematically invalid.''}
Refuting assumptions of independent composition, we establish via
the Welch Simplex bound that homogeneous text probes face a strict
dimensionality ceiling ($D_{\mathrm{eff}} = 1.603$,
$k^* \leq 3.125$), proving that exponential joint-space shrinkage
requires cross-mechanism diversity. The framework's claim is the
bound, not the super-additivity it constrains.

\paragraph{Critique 3: ``An energy ordering inversion on SOTA
models invalidates your core EBM thesis.''}
Evaluations on SOTA outputs reveal an energy-ordering inversion
($\overline{E}_{\mathrm{correct}} > \overline{E}_{\mathrm{incorrect}}$),
suggesting that highly-optimized base models exhibit ``reward
hacking'' dynamics that confound standard linear EBM evaluators.
The in-distribution $\mathrm{AUROC} = 0.9545$ on FoVer is unchanged;
the inversion is an OOD boundary on the energy class, not a
falsification of the framework.

\end{document}
